\section{Related work}
\label{sec:related_work}

We situate our work within four areas. Runtime tools catch
content-level violations but miss structural defects. General-purpose
model checkers could verify the same properties but require manual
modeling. Temporal monitoring is well-studied but not applied to agent
workflows. Agent safety research focuses on alignment, not
orchestration topology. Agentproof fills the gap between runtime
content checking and structural verification.

\subsection{Runtime agent safety tools}

Several tools enforce safety at the point of LLM output or tool invocation.
NVIDIA NeMo Guardrails~\citep{nemo_guardrails} interposes a programmable
dialog rail between the LLM and tool calls, supporting topic control and
output filtering.
Guardrails AI~\citep{guardrails_ai} wraps LLM outputs with validators
that check format, toxicity, and factual consistency.
LlamaGuard~\citep{llamaguard} uses a fine-tuned classifier to detect
unsafe content in model outputs.
Rebuff~\citep{rebuff} focuses specifically on prompt injection detection.
These tools operate at runtime and catch violations only when the
offending code path is actually exercised.
Our approach is complementary: we verify structural properties of the
workflow topology \emph{before deployment}, catching defects such as
unreachable exits, dead-end branches, and missing human gates that
runtime tools cannot detect unless the defective path is triggered
during testing.

\subsection{Static analysis and model checking}

Classical model checking~\citep{clarke1999modelchecking} verifies
finite-state systems against temporal specifications using tools such as
SPIN~\citep{holzmann1997spin}, NuSMV~\citep{cimatti2002nusmv}, and
CBMC~\citep{cbmc}. Abstract interpretation~\citep{cousot1977abstract}
provides sound over-approximations for dataflow analysis.

Table~\ref{tab:comparison_mc} provides a qualitative comparison.
The key difference is \emph{modeling effort}: SPIN requires manual
translation of a workflow into Promela, NuSMV into SMV, and CBMC into
annotated C. For the email triage workflow in our corpus (8~nodes,
7~edges), the equivalent Promela model requires ${\sim}60$ lines and
manual specification of the state space, transitions, and LTL
properties (see \texttt{corpus/comparisons/email\_triage.pml} in our
artifact). Agentproof extracts the same model automatically with a
single function call.

\begin{table}[t]
\centering
\small
\caption{Comparison with general-purpose model checkers.}
\label{tab:comparison_mc}
\begin{tabular}{@{}lllllll@{}}
\toprule
\textbf{Tool} & \textbf{Input} & \textbf{Properties} & \textbf{Modeling} & \textbf{Time} & \textbf{Domain} \\
\midrule
SPIN & Promela (manual) & Full LTL & High & Fast (small) & None \\
NuSMV & SMV (manual) & CTL + LTL & High & Fast (small) & None \\
CBMC & C source & Assertions & Medium & Bounded & None \\
\textbf{Agentproof} & Auto-extracted & Safety LTL fragment & \textbf{None} & $O(|V|{+}|E|)$ & \textbf{Agent workflows} \\
\bottomrule
\end{tabular}
\par\smallskip
\raggedright\footnotesize
SPIN, NuSMV, and CBMC require manual translation of agent workflows into their input languages. Agentproof extracts models automatically from framework APIs, eliminating the modeling step entirely.
\end{table}

Our approach applies model-checking ideas to a new domain: agent
workflow graphs extracted from orchestration framework APIs. The
contribution is not the graph algorithms themselves (which are
standard) but (i)~the automatic extraction from heterogeneous
framework representations, (ii)~a domain-specific property language
tailored to agent safety, and (iii)~empirical evidence that agent
workflows contain structural defects catchable by these methods.

\subsection{Temporal logic monitoring}

The runtime verification (RV) community has developed sophisticated
monitoring frameworks.
JavaMOP~\citep{javamop} monitors Java programs against specifications
in multiple formalisms.
Bauer et al.~\citep{bauer2011ltl3} introduce three-valued LTL
monitoring that distinguishes between ``currently satisfied'' and
``permanently satisfied.''
Barringer et al.~\citep{barringer2004eagle} present the EAGLE framework
for rule-based monitoring.
Our temporal monitors are deliberately lightweight: we target the safety
fragment of LTL with DFAs of 2--$O(k)$ states per rule, trading
expressiveness for deployment simplicity. In our evaluation, all 15
practical agent safety policies fall within this fragment, validating the
design choice empirically (Section~\ref{sec:temporal_eval}).

\subsection{Agent architecture and safety}

Recent work on AI agent safety spans multiple dimensions.
Anthropic's Responsible Scaling Policy~\citep{anthropic_rsp} outlines
safety commitments including evaluation thresholds for autonomous
capabilities.
METR~\citep{metr} provides evaluations for autonomous agent
capabilities and risks.
Research on multi-agent coordination safety~\citep{gabriel2024ethics}
examines the risks of agent-to-agent interaction, including unintended
emergent behavior.
These works focus on alignment, capability evaluation, and policy
governance. Our contribution addresses a complementary layer: the
\emph{structural soundness of the orchestration graph itself}. Even a
well-aligned LLM can produce unsafe outcomes if the workflow graph
routes it through an unintended path or bypasses a required approval
step.

\subsection{Business process verification}

The business process management (BPM) community has extensively studied
workflow soundness. Van der Aalst's work on Petri-net-based soundness
checking for workflow nets~\citep{vanderaalst2011bpm} ensures that
every case can complete and no dead transitions exist---properties
closely analogous to our exit-reachability and dead-end checks. BPMN
verification tools~\citep{dijkman2008bpmn} apply similar analyses to
industry-standard process models.
Our work differs in \emph{domain}: agent workflow graphs have typed
nodes (LLM, TOOL, HUMAN) with framework-specific semantics that BPM
tools do not model, and our extraction pipeline targets programmatic API
objects rather than visual process models.

\subsection{Constrained decoding and structured generation}

Constrained decoding enforces output structure at the token level:
Guidance~\citep{guidance2023} and Outlines~\citep{outlines2023} compile
grammars or JSON schemas into token masks that guarantee well-formed
outputs. These techniques operate \emph{within} a single LLM call,
constraining what the model can generate. Agentproof operates at a
different granularity: it constrains the \emph{transitions between}
computation steps (nodes), not the content generated within any single
step. The two approaches are complementary.
